\section{Height}

In this section, we fix a prime number $p$ and a field $k$ with characteristic $p$.
Let $F$ be a commutative formal group law over $k$. In this section,
unless otherwise specified, all formal group laws are assumed to be
commutative.

In this section, we discuss the definition and properties of the height of a formal group law.

\begin{definition}[Height of homomorphism, {\cite[Chap.~4 Definition 7.1]{silverman2009arithmetic}}]
    Let $F, G$ be formal group laws over $k$, and let $f \in R[\![X]\!]$ be a homomorphism from
    $F(X,Y)$ to $G(X,Y)$. Then the \emph{height of homomorphism} $f$, denoted by $\htt(f)$,
    is the largest number $h$ such that
    \[f(X) = g(X ^ {p ^ h}), \]
    for some power series $g(X)$ over $R$. If $f= 0$, then we set $\htt(f) = \infty$.
\end{definition}

\begin{remark}
    Let $m$ be a natural number prime to $p$, then $\htt([m]) = 0$, since
        \[[m](X) = m X + \cdots\]
\end{remark}

\begin{proposition}[{\cite[Chap.~4 Proposition 7.2]{silverman2009arithmetic}}]
    Let $F, G$ be formal group laws over $k$, and let $f$ be a nonzero
    homomorphism from $F$ to $G$ defined over
    $k$. Then
    \begin{enumerate}
        \item If $f'(0) = 0$, then $f(X) = f_1(X ^ p)$ for some $f_1(X) \in k[\![X]\!]$. That is,
            $\htt(f) \ge 1$.
        \item There is a natural number $m$ such that $f(X) = h(X^{p^m})$ such that
            $h'(0) \neq 0$.
        \item Write $f(X) = g(X^ {p^{\htt(f)}})$. Then $g'(0) \neq 0$.
    \end{enumerate}
    \label{prop: equiv_def_height}
\end{proposition}

\begin{proof}
    \textit{Proof of Part (1):} let $\omega_F, \omega_G$ be the normalized invariant differentials of $F$ and $G$,
    respectively. According to Corollary \ref{cor: omega_hom_a}, we have
    \begin{align*}
        0 &= f'(0) \omega_F \\
         & = \omega_G \circ f
    \end{align*}

    If we write $\omega_G(X) = P_G(X) dX$, then we have
    \[ \omega_G \circ f(X) = \omega_G(f(X)) = P_G(f(X)) f'(X) dX. \]

    Then $f'(X) = 0$, since $P_G(X)$ is of the form $(1 + O(X))$.
    It follows that there is $f_1(X)$ such that
    $f(X) = f_1(X^p)$.

    \textit{Proof of Part (2):} if $f'(0) \neq 0$, then take $m$ to be $0$.
    If $f'(0) = 0$, then by the Part (1) of this proposition, we have that
    $f(X) = f_1(X ^ p)$ for some $f_1(X) \in k[\![X]\!]$. If $f_1'(0) = 0$, then
    we just take $m = 1$. If $f_1'(0) = 0$, we denote $F_p = \varphi_* F$,
    where $\varphi$ is an endomorphism of $k$,
    mapping $a \mapsto a ^ p$. Then $F_p$ is also a formal group law. Now we
    claim that $f_1(X)$ is a homomorphism from $F_p$ to G. We prove this claim
    by the following calculation:
    \begin{align*}
        f_1(F_p(X,Y)) & = f_1 (F(X,Y) ^ p)\\
        & = f_1 (X ^ p) \circ F \\
        & = f(X) \circ F \\
        & = f_1 \circ F.
    \end{align*}

    Then again by the Part (1) of this proposition, we have that there is a $f_2(X) \in k[\![X]\!]$,
    such that $f_1(X) = f_2(X^p)$. Recall that the order of a nonzero power series
    is defined to be the least natural number $n$ sucht that the $n$-th coefficient
    of it is nonzero. Notice that the order $f_i$ decline after each process.
    Then repeating this process at most $n$ times where $n$ is the order
    of $f(X)$, we will get a natural number $m$ satisfying the condition.


    \textit{Proof of Part (3):} we first denote $q = p ^ {\htt(f)}$
    and $F_q = \phi_* F$, where $\phi$ is an endomorphism of $k$,
    mapping $a \mapsto a ^ {q}$. Then $F_q$ is also a formal group law. Now we claim that
    $g(X)$ is a homomorphism from $F_q$ to $G$. We prove this claim by the following calculation:
    \begin{align*}
        g(F_q(X,Y)) & = g (F(X,Y) ^ q)\\
        & = g (X ^ q) \circ F \\
        & = f(X) \circ F \\
        & = g \circ F.
    \end{align*}
    The first equation is guaranteed by the characteristic of $k$.
    If $g'(0) \neq 0$, then by the first assertion of the proposition, there is a
    $g_1(X)$ such that $g(X) = g'(X ^ p)$. It follows that $f(X) = g_1(X^{p^{\htt(f) + 1}})$, which
    contradicts the definition of the height of $f$.
\end{proof}

\begin{remark}
    If we write $f(X) = \sum_{n=1}^{\infty} a_n X ^ n$.
    Let $n$ be the least postive integer such that $a_n$ is nonzero, by Part (2) and (3) of the
    Proposition \ref{prop: equiv_def_height},
    $n$ is of the form $p ^ t$ for some natural number $t$.
    Then we have a alternative description of height of formal group
    homomorphism $f(X)$ from $F$ to $G$, that is the height of $f(X)$ can be
    defined to be the this $t$.
    \label{rem: alter_def_height}
\end{remark}

\begin{definition}[Height of formal group]
    The \emph{height of a formal group law} $F$ is defined to be the height of
    the $p$-series $[p]_F(X)$, denoted by $\htt (F)$.
    % More precisely, let $v_n$ be the $p^n$-th coefficient of the $p$-series $[p](X)$. We say
    % $F(X,Y)$ has height $n$ if $v_i = 0$ for all $i < n$ and $v_n \neq 0$ in $R$.
    % In particular, we say $F(X,Y)$ has infinite height
    % if $[p](X) = 0$.
\end{definition}


% \begin{remark}
%     From the remark above, we have $v_0 = p$. Then $F(X,Y)$ has height $\ge 1$ if and only if
%     $p = 0$ in $R$, which is equivalent to saying either $R$ has characteristic $p$ or $R$ is trivial.
%     Also, $F(X,Y)$ has height exactly $0$ if and only if $p$ is invertible in $R$.
% \end{remark}

\begin{example}
    The $p$-series for the multiplicative formal group
    $\GG_m(X,Y) = (1 + X)(1 + Y) - 1$ is $[p](X) = (1 + X)^p - 1 = X ^ p$, so $\GG_m(X,Y)$
    has height exactly $1$. For the additive formal group law $\GG_a(X,Y)$, the $p$-series is
    $[p](X) = p X = 0$, so $\GG_a(X,Y)$ has infinite height.
\end{example}

\begin{proposition}
    Let $F(X,Y)$ and $G(X,Y)$ be isomorphic formal group laws over $k$. Then $F(X,Y)$ and $G(X,Y)$ have
    the same height.
\end{proposition}

\begin{proof}
    Let $g(X)$ be a formal group isomorphism from $F(X, Y)$ to $G(X, Y)$,
    and let $[p]_F(X)$ and $[p]_G(X)$ be the $p$-series for $F(X, Y)$ and $G(X, Y)$,
    respectively. Note that
    \[ [p]_G(X) = (g \circ [p]_F \circ g^{-1})(X). \]
    Therefore, $[p]_F(X)$ and $[p]_G(X)$ have the same order, and the coefficients of
    their lowest-degree non-zero terms are equal. By Remark \ref{rem: alter_def_height},
    it follows that the heights of $F(X, Y)$ and $G(X, Y)$ are the same.
\end{proof}

% \begin{remark}
%     The additive and multiplicative formal groups are isomorphic over a ring $R$ of
%     characteristic $p$.
% \end{remark}

\begin{proposition}
    Let $F, G$ be formal group law over $k$. Assume that
    $\htt(F) \neq \htt(G)$, then there is only zero homomorphism from $F$ to $G$.
\end{proposition}

\begin{proof}
    We prove the proposition by contradiction. Assume that there is a nonzero homomorphism $f$
    from $F$ to $G$. Note that

    \[ f([p]_F(X)) = [p]_G (f(X)). \]

    This is impossible, since the orders of $[p]_F(X)$ and $[p]_G(X)$ are
    different and $f(X) \neq 0$.
\end{proof}

\begin{proposition}[{\cite[Chap.~4 Proposition 7.3]{silverman2009arithmetic}}]
    Let $F,G,H$ be formal group laws over $k$, and let
    \[F \stackrel{f}{\longrightarrow} G \stackrel{g}{\longrightarrow} H \]
    be a chain of homomorphisms defined over $k$. Then
    \[\htt(g \circ f) = \htt(f) + \htt (g). \]

    \label{prop: height_mul}
\end{proposition}

\begin{proof}
    Denote $f (X) = f_1 (X ^ {p ^ {\htt (f)}})$ and $g(X) = g_1 (X ^ {p ^ {\htt (g)}})$. Then
    \[(g\circ f) (X) = g_1 (f_1 (X ^ {p ^ {\htt (f)}}) ^ {\htt (g)})
        = g_1 (\tilde{f}_1(X ^ {\htt(f) + \htt(g)}))\]
    where $\tilde{f}_1 (X)$ is obtained from $f_1(X)$ raising each coefficient of
    $f_1(X)$ to the $p^{\htt(g)}$-th power.
\end{proof}

\begin{lemma}
    Let $R$ be a commutative ring, $I, J$ be an index set and
    $F \in R[\![X_I]\!]$. Let $f_i, i\in I$ be a set of power series in
    $R[\![X_J]\!]$ with the property that the constant
    coefficient of each $f_i, i\in I$ is nilpotent in $R$.
    Let $G(X_J)$ be the power series obtained by substituting
    $f_i$ into $F$. Then we have that
    the order of $G(X_J)$ is at least the multiplication of the order of $F$
    and the least of the order of $f_i$.

    \label{lem: order_subst}
\end{lemma}

\begin{proof}
    Recall that for any power series $f(X)$, the order of $f(X)$ is at least
    $n$ is equivalent to for all $j < n$, the $j$-th homogeneous component
    $f(X)$ is zero.

    Then it is enough to show that for any $j$ less than the multiplication of
    the order of $F$ and the least of the order of $f_i$, then the $j$-th homogeneous
    component of $G$ is zero. However, the least possible nonzero homogeneous component in
    $G(X)$ is at least of degree the multiplication of
    the order of $F$ and the least of the order of $f_i$.
\end{proof}

\begin{proposition}[{\cite{hazewinkel1978formal}, Definition 18.3.2}]
    Let $F$ be a formal group law over $k$, $f$ and $g$ be two formal group
    endomorphisms of $F$. Then $F(f, g)$ is an endomorphism of $F$ and
    $\htt (F(f,g)) \ge \min \{\htt (f), \htt(g)\}$.
    \label{prop: height_add}
\end{proposition}

\begin{proof}
    First, we show that $F(f,g)$ is an endomorphism of $F$. It follows from
    the following calculation:
    \begin{align*}
        F(f(X), g(X)) \circ F(X,Y) &= F(f(F(X,Y)), g(F(X,Y))) \\
        & = F(F(f(X),f(Y)), F(g(X), g(Y))) \\
        & = F(F(f(X), g(X)), F(f(Y), g(Y))) \\
        & = F(X,Y) \circ F(f(X), g(X)).
    \end{align*}
    By Lemma \ref{lem: order_subst} the fact that order of $F(X,Y)$ is $1$,
    then the order $F(f(X), g(X))$ is lower order $f$ and $g$. Therefore,
    by Remark \ref{rem: alter_def_height}, it follows that
    $\htt (F(f,g)) \ge \min \{\htt (f), \htt(g)\}$.
\end{proof}

\begin{remark}
    Let $i(X)$ be the additive inverse of formal group law $F$, i.e. the power
    series such that $F(X, i(X)) = 0$.
    Note that
    \[i(X) \circ F(X,Y) = i(F(X,Y)) = F(i(Y), i(X)) = F(i(X), i(Y)). \]

    Then $i(X)$ is an endomorphism of $F$. By Proposition \ref{prop: height_add} and
    \ref{prop: height_mul}, we have that set of all endomorphisms over $k$ of $F$ form a
    ring, denote as $\End_k(F)$. The multiplication is given by composition of power series and the addition
    of $f, g \in \End_k(F)$ is defined to be $F(f,g)$. And the height function
    $\htt$ defines a valuation on the ring $\End_k(F)$.
\end{remark}

Let $k$ be an algebraic closed field with characteristic $p$,
formal group laws over $k$ is classified by the heights of them.

\begin{theorem}[Lazard {\cite[Theorem IV]{lazard1955sur}}]
    Let $k$ be an algebraically closed field of characteristic $p$. Then two formal group laws
    $F(X,Y), G(X,Y)$ over $k$ are isomorphic if and only if they have the same height.
\end{theorem}


